% 大熊座（综述）
% license CCBYSA3
% type Wiki

本文根据 CC-BY-SA 协议转载翻译自维基百科\href{https://en.wikipedia.org/wiki/Ursa_Major}{相关文章}。

大熊座（Ursa Major），亦称“巨熊座”，是位于北天的一座星座，其相关神话很可能可以追溯至史前时代。其拉丁名称意为“较大的（或更大的）熊”，用以指代并区别于附近的小熊座（Ursa Minor）。$^{[2]}$ 在古代，大熊座是公元2世纪托勒密所列出的最初48个星座之一，其编制基础源自更早的希腊、埃及、巴比伦与亚述天文学家的研究成果。$^{[3]}$在现代被正式划定的88个星座中，大熊座按面积排名第三。

大熊座最为人熟知的是由其七颗主星构成的星群，这一星群在不同文化中被称为“大斗”“马车”“查理的车”或“犁”等名称。尤其值得注意的是，“大斗”的恒星排列形态在视觉上与“小斗”相呼应。其中的两颗恒星——天枢与天璇（分别为大熊座$\alpha$星与$\beta$星，Dubhe 与 Merak）——可作为导航指示星，用来指向当前的北极星，即位于小熊座中的北极星（Polaris）。

大熊座及其所包含或交叠的星群在世界诸多文化中具有重要意义，常被视为“北方”的象征。阿拉斯加州旗帜上对大熊座的描绘，便是这一象征在现代的一个实例。

在北半球的大多数地区，大熊座全年可见，并在中北纬度地区呈现为拱极星座。从南半球的温带纬度看，其主要星群不可见，但星座的南部区域仍可被观测到。
\subsection{特征}  
大熊座覆盖约1279.66平方度，约占全天空的3.10\%，是第三大星座。$^{[4]}$ 1930年，欧仁·德尔波特为其划定了国际天文学联合会（IAU）认可的正式星座边界，将其定义为一个由28条边组成的不规则多边形。在赤道坐标系中，该星座的赤经范围介于$08^{\mathrm h}08.3^{\mathrm m}$至$14^{\mathrm h}29.0^{\mathrm m}$之间，赤纬范围介于$+28.30^\circ$至$+73.14^\circ$ 之间。$^{[5]}$大熊座与八个星座相邻：北方与东北方为天龙座，东方为牧夫座，东与东南为猎犬座，东南为后发座，南方为狮子座与小狮座，西南为天猫座，西北为鹿豹座。国际天文学联合会于1922年采用其三字母缩写“UMa”。$^{[5]}$
\subsection{特征}  
\subsubsection{星群}
\begin{figure}[ht]
\centering
\includegraphics[width=6cm]{./figures/f2d88b83474133e7.png}
\caption{大熊座与北极星，以及北斗七星中明亮恒星的名称} \label{fig_UrMa_1}
\end{figure}
\begin{figure}[ht]
\centering
\includegraphics[width=6cm]{./figures/3eb25f5a42cf8f25.png}
\caption{肉眼所见的大熊座星座} \label{fig_UrMa_2}
\end{figure}
大熊座中七颗明亮恒星的轮廓构成了一个星群，在美国和加拿大被称为“北斗”（Big Dipper），而在英国则被称为“犁”（Plough）$^{[6]}$，或在历史上称为“查理之车”（Charles' Wain）$^{[7]}$。这七颗恒星中有六颗的视星等为二等或更亮，使其成为夜空中最为人熟知的恒星图案之一。$^{[8][9]}$正如这些常见名称所暗示的那样，其形状常被认为类似于一把勺子、一种农用犁具，或一辆马车。在大熊座的语境中，它们通常被绘制为巨熊的后躯与尾巴。从“勺斗”部分开始，沿顺时针方向（即天空中的东方方向）延伸至柄部，这些恒星依次为：
\begin{itemize}
\item 天枢（Dubhe，“熊”），视星等为1.79，是全天第35亮的恒星，也是大熊座中第二亮的恒星。  
\item 天璇（Merak，“熊的腰部”），视星等为2.37。  
\item 天玑（Phecda，“大腿”），视星等为2.44。  
\item 天权（Megrez），其名称意为“尾巴的根部”，指其位于熊的躯体与尾巴的交汇处（或北斗中斗身与柄部的连接点）。  
\item 玉衡（Alioth），其名称并非指熊，而是意为“黑马”，这是由原名讹变并误用于与其名称相近的伴星开阳（Alcor）所致，后者是肉眼可见的双星伴侣，围绕着开阳主星运行。$^{[10]}$ 玉衡是大熊座中最亮的恒星，也是全天第33亮的恒星，视星等为1.76。它同时也是化学丰度异常的 Ap 型恒星中最明亮的一颗，这类磁性恒星的某些化学元素会出现异常贫化或富集，并且随着恒星自转而表现出变化。$^{[10]}$
\item 开阳（Mizar，大熊座$\zeta$星），位于北斗柄端向内数第二颗，是大熊座中第四亮的恒星。“Mizar”一名意为“腰带”。它与其视伴星辅星（Alcor，大熊座80号星）共同构成著名的双星系统，在阿拉伯传统中被称为“马与骑士”。  
\item 摇光（Alkaid，大熊座$\eta$星），位于熊尾的末端。其视星等为1.85，是大熊座中第三亮的恒星。$^{[11][12]}$
\end{itemize}
除天枢与摇光之外，北斗七星中的其他恒星都具有指向人马座方向的共同自行运动。人们还识别出少数具有相同运动特征的其他恒星，它们共同被称为“大熊座移动星群”（Ursa Major Moving Group）。
\begin{figure}[ht]
\centering
\includegraphics[width=6cm]{./figures/58bc21206ce95031.png}
\caption{大熊座与小熊座相对于北极星的位置关系} \label{fig_UrMa_3}
\end{figure}
天璇（Merak，大熊座$\beta$星）与天枢（Dubhe，大熊座$\alpha$星）被称为“指极星”，因为它们有助于寻找北极星（亦称北辰星或极星）。通过目视从天璇指向天枢画一条直线（记为1个单位），并沿同一方向继续延伸5个单位，视线便会落在北极星上，从而准确指示真北方向。

在阿拉伯文化中，还识别出另一组星群，被视为一只跳跃羚羊留下的三对足迹。$^{[13]}$ 这是一列位于星座南部边界的三对恒星。从东南至西南依次为：“第一次跳跃”，由$\nu$与$\xi$大熊座（分别为 Alula Borealis 与 Alula Australis）组成；“第二次跳跃”，由 $\lambda$与$\mu$ 大熊座（Tania Borealis 与 Tania Australis）组成；以及“第三次跳跃”，由 $\iota$与$\kappa$ 大熊座（Talitha 与 Alkaphrah）组成。同一星群在中文中称为“三台”$^{[13]}$，在印度传统中称为 Trivikrama$^{[14]}$，两者均意为“三步”。
\subsubsection{其他恒星}  
W 大熊座是接触双星变星的一类原型，其视星等在 7.75 与 8.48 之间变化。

47 大熊座是一颗类太阳恒星，拥有一个由三颗行星组成的系统。$^{[15]}$其中，47 大熊座 b 于1996年发现，其公转周期为1078天，质量为木星的2.53倍。$^{[16]}$47 大熊座 c 于2001年发现，其公转周期为2391天，质量为木星的0.54倍。$^{[17]}$47 大熊座 d 于2010年发现，其周期尚不确定，介于8907至19097天之间，质量为木星的1.64倍。$^{[18]}$该恒星的视星等为5.0，距离地球约46光年。$^{[15]}$

恒星 TYC~3429-697-1（$9^{\mathrm h}40^{\mathrm m}44^{\mathrm s}$，$48^\circ14'2''$），位于 $\theta$ 大熊座以东、北斗以西南的位置，被认定为美国特拉华州的州星，并被非正式地称为“特拉华钻石”。$^{[19]}$
\subsubsection{深空天体}
\begin{figure}[ht]
\centering
\includegraphics[width=6cm]{./figures/a11c8340dbe9123a.png}
\caption{旋涡星系} \label{fig_UrMa_4}
\end{figure}
大熊座中分布着多座明亮的星系，其中包括位于熊头上方的一对星系——梅西耶81（M81，夜空中最明亮的星系之一）与梅西耶82（M82），以及位于摇光（Alkaid）东北方向的旋涡星系（M101）。该星座中还包含旋涡星系梅西耶108（M108）与梅西耶109（M109）。明亮的行星状星云——猫头鹰星云（M97）则位于北斗斗身的底部附近。

M81 是一座几乎正面对着地球的旋涡星系，距离地球约1180万光年。与大多数旋涡星系类似，它的核心由老年恒星构成，而旋臂中则充满年轻恒星与星云。M81 与 M82 同属距离本星系群最近的一个星系团。

M82 是一座几乎侧向可见的星系，正与 M81 发生引力相互作用。它是夜空中最明亮的红外星系。$^{[20]}$2014年1月21日，人们在 M82 中观测到了一次视星等显著的 Ia 型超新星 SN~2014J。$^{[21]}$

M97，又称猫头鹰星云，是一座距离地球约1630光年的行星状星云，其视星等约为10。它于1781年由皮埃尔·梅尚发现。$^{[22]}$

M101，又称旋涡星系，是一座正面对地球的旋涡星系，距离约2500万光年，由皮埃尔·梅尚于1781年发现。其旋臂中存在大规模恒星形成区域，并伴随强烈的紫外辐射。$^{[20]}$该星系的综合视星等为7.5，可通过双筒望远镜或天文望远镜观测，但肉眼不可见。$^{[23]}$

NGC~2787 是一座距离约2400万光年的透镜状星系。与多数同类星系不同的是，它在中心具有明显的棒状结构，并且拥有一个由球状星团构成的晕，表明其具有较高的年龄与相对稳定性。$^{[20]}$

NGC~2950 是一座透镜状星系，距离地球约6000万光年。

NGC~3000 是一对双星，但在星表中被编入星云型天体。

NGC~3079 是一座位于约5200万光年之外的星暴旋涡星系。其中心存在一个马蹄形结构，暗示着超大质量黑洞的存在，该结构由黑洞产生的超风所形成。$^{[20]}$

NGC~3310 是另一座星暴旋涡星系，距离地球约5000万光年。其明亮的白色外观源于异常旺盛的恒星形成速率，这一过程始于约一亿年前的一次星系并合。对该星系及其他星暴星系的研究表明，其星暴阶段可持续数亿年，远长于以往的假设。$^{[20]}$

NGC~4013 是一座侧向可见的旋涡星系，距离地球约5500万光年。它具有显著的尘埃带，并包含多个清晰可见的恒星形成区域。$^{[20]}$

I~Zwicky~18 是一座年轻的矮星系，距离地球约4500万光年。作为可见宇宙中已知最年轻的星系，I~Zwicky~18 的年龄约为400万年，仅为太阳系年龄的约千分之一。它内部充满恒星形成区，以极高的速率生成大量炽热、年轻的蓝色恒星。$^{[20]}$

哈勃深场（Hubble Deep Field）位于大熊座$\delta$星的东北方向。
\subsubsection{流星雨}  
$\alpha$大熊座流星雨（Alpha Ursae Majorids）是该星座中的一场较小型流星雨。$^{[24]}$它们可能由彗星 C/1992~W1（Ohshita）所引发。$^{[24][25]}$

$\kappa$ 大熊座流星雨（Kappa Ursae Majorids）是一场新近发现的流星雨，其极大期位于11月1日至11月10日之间。$^{[26]}$

10月大熊座流星雨（October Ursae Majorids）于2006年由日本研究人员发现。它们可能由一颗长周期彗星所致。$^{[27]}$该流星雨的极大期介于10月12日至10月19日之间。$^{[28]}$
\subsubsection{系外行星}  
HD~80606 是一颗类太阳恒星，位于一个双星系统中，与其伴星 HD~80607 围绕共同的质心运行，两者的平均间距约为1200~AU。2003年的研究表明，其唯一已知行星 HD~80606~b 是一颗未来的“热木星”，模型显示它最初在距离其母星约5~AU、且轨道几乎垂直的轨道上演化而来。该行星质量约为木星的4倍，预计最终会在 Kozai 机制作用下演化为一条更加圆形、且与恒星自转轴更为对齐的轨道。然而在当前阶段，它仍处于一条极端高偏心率的轨道上，其远日点距离约为1个天文单位，而近日点则仅约为6个恒星半径。$^{[29]}$
\subsection{历史}
\begin{figure}[ht]
\centering
\includegraphics[width=6cm]{./figures/e4987f34375d0ea2.png}
\caption{约1700年刻于石上的大熊座，地点为苏格兰法夫郡克雷尔（Crail）} \label{fig_UrMa_5}
\end{figure}
大熊座已被重建为一个印欧语系的星座。$^{[30]}$它是公元2世纪天文学家托勒密在其《天文学大成》（Almagest）中所列出的48个星座之一，托勒密将其称为 Arktos Megale（“大熊”）。$^{[a]}$
\subsection{神话}
大熊座这一星座被多个彼此关联的文明视为一只熊，通常是雌熊。$^{[32][33]}$这可能源于一个共同的口述传统，即“宇宙狩猎”神话，其历史可追溯至一万三千年以上。$^{[34]}$朱利安·迪伊（Julien d'Huy）利用统计学与系统发生学工具，重建了该故事的旧石器时代版本：“有一种动物，是一种有角的食草动物，尤其是麋鹿。一名人类追逐这只偶蹄动物。狩猎通向天空或发生在天空中。当该动物被转化为星座时，它仍然是活着的。它形成了北斗七星。”$^{[35]}$
\subsubsection{阿拉伯民间传统}  
尽管伊斯兰教出现之前的阿拉伯人将更大的星座——大熊座——识别为一只熊，这或许受到了希腊文化的影响，但他们在传统上始终将北斗七星与小熊座视为一对对应之物。二者都被想象为送葬队伍，其中“斗”的部分被视为灵柩，而“柄”则是随行的哀悼者队列。北斗七星被称为 banāt an-naʿsh al-kubrā，字面意为“灵柩的年长女儿们”。然而这里的“女儿”指的是隶属于灵柩之物，即哀悼者，因此更恰当的译法是“较大的送葬队伍”，而小熊座则被称为“较小的送葬队伍”。另有一种传说认为，灵柩上躺着的是队伍中众人的父亲，他名叫 Naʿash，被 Al-judayy（北极星的阿拉伯名称）所杀，而这支送葬队伍正是在追逐凶手。$^{[36]}$
\subsubsection{希腊—罗马传统}  
在希腊神话中，宙斯（诸神之王，在罗马神话中称为朱庇特）爱慕一位名叫卡利斯托（Callisto）的年轻女子，她是阿耳忒弥斯的宁芙（罗马人称其为狄安娜）。宙斯嫉妒的妻子赫拉（罗马神话中的朱诺）发现卡利斯托因被宙斯强暴而生下了儿子阿耳卡斯（Arcas），于是将卡利斯托变成一只熊作为惩罚。$^{[37]}$处于熊形态的卡利斯托后来与她的儿子阿耳卡斯相遇，阿耳卡斯几乎将这只熊刺杀。为了避免悲剧发生，宙斯将二人一同送上天空，使卡利斯托化为大熊座，而阿耳卡斯化为牧夫座。奥维德称大熊座为“帕拉西亚之熊”（Parrhasian Bear），因为卡利斯托来自阿卡迪亚的帕拉西亚地区，该神话正发生于此。$^{[38]}$

希腊诗人阿拉托斯将该星座称为 Helikē（意为“旋转”或“缠绕”），因为它围绕天极旋转。《奥德赛》中指出，它是唯一一个永不沉入地平线之下、并且“沐浴在大洋波涛之中”的星座，因此常被用作导航的天球参照物。$^{[39]}$它还被称为“车”（Wain）或“Plaustrum”，这是一个拉丁词，指由马牵引的货车。$^{[40]}$
\subsubsection{印度教传统}  
在印度教中，对大熊座／北斗七星／巨熊座的最早记载被称为“七仙人”（Saptarshi），其中每一颗恒星分别代表一位七仙人（Rishi），即：婆利古（Bhrigu）、阿特里（Atri）、安吉拉斯（Angiras）、婆悉吒（Vasishtha）、普拉斯提亚（Pulastya）、普拉哈（Pulaha）与克拉图（Kratu）。这一记载见于《梨俱吠陀》（约公元前1500—1200年），该文献是人类已知最古老的文本之一。

古印度文献中的相关记载包括：
\begin{enumerate}
\item 《梨俱吠陀》（第一曼荼罗，第24赞歌，第10节），在宇宙秩序与天体意义的语境中提及七仙人。
\item 《摩诃婆罗多》（约公元前400年至公元400年），论及“七仙人星群”（Saptarishi Mandal）作为航行与指引的恒星。
\item 《往世书》（如《毗湿奴往世书》《梵天卵往世书》等），将七仙人描述为掌握宇宙智慧的神圣贤者。
\end{enumerate}
至于小熊座，其并未在早期吠陀文献中被明确提及，但在后来的天文学著作中得到了确认，例如：
\begin{enumerate}
\item 《吠檀伽天文书》（Vedanga Jyotisha，约公元前1400—1200年）
\item 《苏利耶悉檀多》（Surya Siddhanta，约公元4世纪）
\end{enumerate}
关于大熊座前方两颗恒星指向北极星的事实，被解释为毗湿奴赐予少年圣者德鲁瓦（Dhruva）的恩典。$^{[41]}$

因此，《梨俱吠陀》构成了关于大熊座的最早文字记载，而小熊座则是在后期天文学传统中逐渐获得重要地位。
\subsubsection{在犹太教与基督教中}  
大熊座可能在《圣经·约伯记》中被提及，该书成书时间一般认为介于公元前7世纪至公元前4世纪之间，但这一对应关系常存在争议。$^{[42]}$
\subsubsection{东亚传统}  
在中国与日本，北斗七星被称为“北斗”（中文：běidǒu，日文：hokuto）。在古代，七颗恒星各自都有专名，多源自中国古代星名体系：

\begin{itemize}
\item “枢”（樞，中文：shū；日文：sū）对应天枢（Dubhe，大熊座$\alpha$星）  
\item “璇”（璇，中文：xuán；日文：sen）对应天璇（Merak，大熊座$\beta$星）  
\item “玑”（璣，中文：jī；日文：ki）对应天玑（Phecda，大熊座$\gamma$星）  
\item “权”$^{[43]}$（權，中文：quán；日文：ken）对应天权（Megrez，大熊座 $\delta$ 星）  
\item “玉衡”（中文：yùhéng；日文：gyokkō）对应玉衡（Alioth，大熊座$\epsilon$ 星）  
\item “开阳”（開陽，中文：kāiyáng；日文：kaiyō）对应开阳（Mizar，大熊座$\zeta$ 星）  
\item 摇光（Alkaid，大熊座$\eta$星）则拥有多个别名：“剑”（劍，中文：jiàn；日文：ken，为“剑先”劍先的简化形式）、“摇光”（搖光，中文：yáoguāng；日文：yōkō），以及“破军星”（破軍星，中文：pójūn xīng；日文：hagun sei），因为在传统观念中，军队若向该星方向行军被视为不祥之兆。$^{[44]}$
\end{itemize}
在神道教中，大熊座中最亮的七颗恒星隶属于天之御中主神（Ame-no-Minakanushi），被视为最古老且最具力量的神祇。

在韩国，该星座被称为“北方七星”。相关神话讲述了一位寡妇与其七个儿子的故事：寡妇爱上一位鳏夫，但通往其住所需横渡一条河流。七个儿子体恤母亲，在河中放置踏石。母亲并不知情，却为这些踏石祝福；待七个儿子去世后，他们化为天上的星座。
\subsubsection{美洲原住民传统}  
易洛魁人将玉衡（Alioth）、开阳（Mizar）与摇光（Alkaid）视为三名猎人，正在追逐巨熊。按照其中一种神话版本，第一位猎人（玉衡）手持弓箭，准备射杀巨熊；第二位猎人（开阳）肩扛一口大锅——即其伴星辅星（Alcor）——用于烹煮猎物；第三位猎人（摇光）则拖拽着一捆柴火，用以在锅下生火。

拉科塔人称该星座为 Wičhákhiyuhapi，意为“巨熊”。

根据托马斯·莫顿在《新英格兰迦南》中的记载，万帕诺亚格人（阿尔冈昆语族）将大熊座称为 “maske”，意为“熊”。$^{[45]}$

瓦斯科—威什拉姆族原住民则将该星座解释为五只狼与两只熊，它们被郊狼神留在了天空中。$^{[46]}$
\subsubsection{日耳曼传统}  
在北欧异教传统中，北斗七星被称为 Óðins vagn，即“奥丁之车”。同样地，奥丁在诗性称谓（Kennings）中也被称为 vagna verr（“车之守护者”）或 vagna rúni（“车之密友”）。$^{[47]}$
\subsubsection{乌拉尔语系传统}  
在芬兰语中，这一星群有时以其古老的芬兰名称 Otava 被称呼。该名称在现代芬兰语中的含义几乎已经被遗忘，其原意是“鲑鱼拦网”。古代芬兰人认为，熊（Ursus arctos）是被装在金篮中从大熊座降到地上的；当一头熊被猎杀后，人们会将熊头放置在树上，以便让熊的灵魂返回大熊座。

在北欧的萨米语中，这一星座的一部分（即北斗七星去除天枢与天璇）被视为伟大猎人 Fávdna 的弓（猎人对应的恒星为大角星 Arcturus）。在主要的萨米语言——北萨米语中，它被称为 Fávdnadávgi（“Fávdna 之弓”），或简称为 dávggát（“弓”）。这一星座在萨米族国歌中占据重要地位，歌词开头为 Guhkkin davvin dávggaid vuolde sabmá suolggai Sámieanan，其意为“在遥远的北方，在弓之下，萨米之地缓缓显现”。“弓”在萨米族关于夜空的传统叙事中具有核心意义，其中多位猎人试图追逐 Sarva——巨大的驯鹿星座，它几乎占据了半个天空。根据传说，Fávdna 每晚都准备拉弓射击，但他迟疑不决，因为他可能会击中被称为 Boahji（“铆钉”）的北极星（Stella Polaris），这将导致天空崩塌并终结世界。$^{[48]}$
\subsubsection{东南亚传统}  
在缅甸语中，Pucwan Tārā（ပုဇွန် တာရာ，[bəzʊ̀ɴ tàjà]）是一个由大熊座头部与前肢恒星组成的星座名称；pucwan（ပုဇွန်）是对甲壳类动物的统称，如对虾、虾、蟹、龙虾等。

在爪哇语中，该星座被称为 “lintang jong”，意为“舢板星座”；在马来语中则被称为 “bintang jong”。$^{[49]}$
\subsubsection{神秘学传统}  
在神智学（Theosophy）中，人们相信昴宿星团的七颗星将来自银河洛格斯（Galactic Logos）的七道灵性能量汇聚并传导至大熊座的七颗恒星，再传至天狼星，随后传至太阳，再到地球之神（Sanat Kumara），最终通过七道光线的七位大师传递给人类。$^{[50]}$
\subsection{在文化中}  
大熊座曾被荷马、斯宾塞、莎士比亚、丁尼生等诗人提及，也出现在费德里科·加西亚·洛尔迦的《献给月亮之歌》中。$^{[51]}$古代芬兰诗歌同样提到这一星座，它还出现在文森特·梵高的绘画作品《罗讷河上的星夜》中。$^{[52][53]}$
\subsubsection{图形可视化}
